This section reports empirical results under the evaluation protocol described above.
All experiments share the same data splits, model family, entity set, and training protocol.

Time aggregation yields a clear lift over a target-encoding-only baseline under the no-lookahead protocol.
Appendix Figure \ref{fig:te_lift_deltas} reports paired deltas with confidence intervals for this comparison.
Table \ref{tab:absolute_metrics_selected} reports the corresponding absolute test metrics for selected specifications.
Appendix Table \ref{tab:top_specs_absolute} reports additional absolute metrics for representative high-performing specifications.

\begin{table}[t]
  \centering
  \small
  \resizebox{\linewidth}{!}{%
  \begin{tabular}{lcccccc}
    \toprule
    & \multicolumn{3}{c}{Test ROC AUC} & \multicolumn{3}{c}{Test PR AUC} \\
    \cmidrule(lr){2-4} \cmidrule(lr){5-7}
    Specification & A & B & mean & A & B & mean \\
    \midrule
    TE only &
      0.74175 & 0.74382 & 0.74278 &
      0.36931 & 0.34796 & 0.35863 \\
    TE plus trailing $(1, 6, 24, 48, 168)$ &
      0.74831 & 0.75198 & 0.75014 &
      0.37875 & 0.35629 & 0.36752 \\
    TE plus trailing $(1, 6, 24, 48, 168)$ plus event50 &
      0.74864 & 0.75245 & 0.75054 &
      0.37844 & 0.35745 & 0.36794 \\
    \bottomrule
  \end{tabular}
  }
  \caption{Absolute test ROC AUC and PR AUC on Fold A and Fold B for selected specifications under the TE plus time aggregation setting.}
  \label{tab:absolute_metrics_selected}
\end{table}

\subsection{Window design effects without target encoding}
Figure \ref{fig:trimmed_note_rocauc} provides a compact view of window design effects in the setting without target encoding.
It reports test ROC AUC averaged over folds, and it marginalizes over the other design dimension.

\begin{figure}[t]
  \centering
  \includegraphics[width=\linewidth]{figures/fig_trimmed_note_test_rocauc_by_length_tuple.pdf}
  \vspace{0.5em}
  \includegraphics[width=\linewidth]{figures/fig_trimmed_note_test_rocauc_by_shape.pdf}
  \caption{No target encoding sensitivity view. Panel A reports test ROC AUC by window length tuple, averaged over shapes and folds. Panel B reports test ROC AUC by window shape, averaged over length tuples and folds.}
  \label{fig:trimmed_note_rocauc}
\end{figure}

\subsection{Window design summary}
Figure \ref{fig:traffic_light} summarizes shape effects relative to trailing windows within each length tuple.
Event count windows are positive in most settings.
Gap and bucket shapes are negative in most settings.
Calendar alignment effects are small and mixed.

\begin{figure}[t]
  \centering
  \includegraphics[width=\linewidth]{figures/fig_traffic_light_shape_contrasts.pdf}
  \caption{Traffic light summary of shape effects relative to trailing windows, split by whether target encoding is used.}
  \label{fig:traffic_light}
\end{figure}

\subsection{Window design choices with target encoding}
The next comparison evaluates window design choices with target encoding enabled under the same evaluation protocol.
Trailing windows with length tuple $(1, 6, 24, 48, 168)$ are a strong baseline.
Event count windows are the only consistently helpful shape modification.
Gap windows and bucketized windows underperform trailing windows within the same length tuple.

Figure \ref{fig:league_table} provides a ranked view of best and worst specifications using paired deltas in test ROC AUC.

\begin{figure}[t]
  \centering
  \includegraphics[width=\linewidth]{figures/fig_league_table_top_bottom.pdf}
  \caption{League table of best and worst specifications by paired deltas in test ROC AUC relative to trailing $(1, 6, 24, 48, 168)$.}
  \label{fig:league_table}
\end{figure}

Appendix Figure \ref{fig:decision_plot} provides a detailed view that separates the effect of length tuple choices from the effect of shape choices.
It reports paired deltas and fold specific estimates.

\subsection{Interaction with target encoding}
Target encoding yields a large uplift relative to time aggregation only models in this setup.
Appendix Figure \ref{fig:te_uplift} reports the uplift across specifications.
Even with target encoding, event count windows provide a small additional gain in ROC AUC over trailing windows.
The gain is $0.00033$ on Fold A and $0.00047$ on Fold B in the paired tests.
PR AUC gains from event count windows are small and mixed.

\subsection{Practical guidance}
For this dataset and protocol, a trailing window design is a robust default.
An event count window can be added when a small additional ROC AUC gain is valuable.
Gap windows and bucketized windows are not recommended under no-lookahead evaluation.

\clearpage
